\section{What the Objective Rules Out, and What It Does Not}
\label{sec:coverage}

An earlier version of this work argued that reverse-reachability sampling has no domain here
because every RR set is empty.  That argument is withdrawn.  The probe it rested on required a
predecessor's window to contain zero and excluded the root, and it tested whether the returned set
was non-empty rather than whether it intersected the seed set.  A reverse-reachability set for a
root answers \emph{which seeds can reach this root}; it does not require a node to self-activate
under zero incoming influence.  The empty sets were an artefact of a mis-specified probe.

What survives is a strictly narrower statement, proved by exact arithmetic on a one-hop instance
rather than by simulation.

\subsection{The instance}

Take a single target $D = \{t\}$ and two candidates $a, b$ with edges $a \to t$ and $b \to t$,
both of weight $\omega = 0.40$, and seeds drawn from $V \setminus D$.

Under the model's own window distribution -- $(\theta^\kappa, \theta^\tau)$ uniform on the
2-simplex $\{0 \le \theta^\kappa \le \theta^\tau \le 1\}$ -- the target's positive-activation
probability is $g(\delta) = 2\delta(1-\delta)$ with $\delta$ the exposure it receives, so

\begin{center}
\begin{tabular}{lccc}
\toprule
$S$ & $\delta$ & $F_D(S)$ under the model & $F_D(S)$ under uniform-threshold LT \\
\midrule
$\emptyset$   & $0.00$ & $0.0000$ & $0.0000$ \\
$\{a\}$       & $0.40$ & $0.4800$ & $0.4000$ \\
$\{a,b\}$     & $0.80$ & $0.3200$ & $0.8000$ \\
\bottomrule
\end{tabular}
\end{center}

Two things follow, and they must be kept apart.

\textbf{The objective is non-monotone.} $F_D(\{a,b\}) = 0.3200 < F_D(\{a\}) = 0.4800$.
Adding a seed can strictly decrease the target's activation probability.

\textbf{No non-negative seed-independent coverage function matches it.} For any random set $R$
independent of $S$, the function $S \mapsto c \cdot \Pr[S \cap R \neq \emptyset]$ is non-decreasing
and submodular in $S$, because adding a seed can only create intersections and the marginal value
of a seed cannot increase as $S$ grows.  A non-monotone $F_D$ therefore cannot equal such a
function on this instance, for any $c$ and any distribution of $R$.

\subsection{Scope}

The statement above excludes exactly one thing: the standard, non-negative, seed-independent
coverage representation that underlies RR sampling.  It does not exclude, and should not be read
as excluding:

\begin{itemize}
  \item signed decompositions, or sampling structures other than non-negative coverage;
  \item algorithms specialised to restricted instance classes;
  \item Monte-Carlo estimation.  A plain sample mean is an unbiased estimator of $F_D$ and is the
        reference used throughout this paper;
  \item any claim that a learned predictor must help.  Whether learning adds anything is an
        empirical question this paper treats as open.
\end{itemize}

\subsection{Why the threshold distribution must be stated explicitly}

The table also shows why the two threshold readings cannot be conflated.  Under
uniform-threshold LT with $\theta^\kappa \sim U[0,1]$ and no upper threshold, the same instance is
\emph{monotone} ($0.4000 \to 0.8000$), so the argument would fail there.  The
source model samples $(\theta^\kappa, \theta^\tau)$ jointly from the simplex, which induces the
marginal $F^\kappa(x) = 2x - x^2$ for the lower threshold and $F^\tau(x) = x^2$ for the upper; a
later lemma in the same work describes sampling $\theta^\tau$ uniformly and then
$\theta^\kappa$ conditionally, which is a different joint law.  Every claim in this paper
therefore names the distribution it assumes, and the degenerate no-overexposure setting is
treated as its own configuration rather than as standard uniform-threshold LT.

The practical consequence for the rest of the paper is that this objective must be estimated by
simulating the process.  Sampling-based coverage machinery is unavailable, so the question turns
into how to spend a simulation budget well -- which is the subject of \S\ref{sec:experiments}.
